\documentclass[border=1pt]{standalone}
\usepackage{mathptmx}  % Times text and math
\usepackage{circuitikz}
\ctikzset{bipoles/length=0.75cm, transistors/scale=0.9,
  inductors/coils=4, current/distance=.5, bipoles/diode/height=.32,
  bipoles/diode/width=.26}
\begin{document}
\fontsize{8}{9.6}\selectfont
% llc_fullbridge: coordinates are layout units, 1 unit = 1.4 cm across and 1.2 cm up (ieee-double)
\begin{circuitikz}[x=1.4cm, y=1.2cm, line width=0.55pt, cute inductors, european resistors, american voltages]
% ---- wires ----
\draw (0,3.9) -- (2.35,3.9);
\draw (0,0) -- (2.35,0);
\draw (1,2.9) -- (1,1);
\draw (2.35,2.9) -- (2.35,1);
\draw (1,2.5) -- (2.2643,2.5);
\draw (2.4357,2.5) -- (2.6,2.5);
\draw (4.5,2.5) -- (5.15,2.5);
\draw (5.15,1.2) -- (2.35,1.2);
\draw (6.55,0) -- (6.55,1.5);
\draw (7.6,2.3) -- (7.6,3.9);
\draw (5.85,2.5) -- (6.55,2.5);
\draw (5.85,1.2) -- (6.4643,1.2);
\draw (6.6357,1.2) -- (7.6,1.2);
\draw (6.55,3.9) -- (9.7,3.9);
\draw (6.55,0) -- (9.7,0);
% ---- Vin (V) ----
\draw (0,3.9) to[V, l_=$V_\mathrm{in}$] (0,0);
% ---- S1 (nmos) ----
\node[nmos,bodydiode] (S1) at (1,3.4) {};
\draw (1,3.9) -- (S1.D) (S1.S) -- (1,2.9);
\node[anchor=west] at ([xshift=0.33cm]S1.center) {$S_{1}$};
% ---- S2 (nmos) ----
\node[nmos,bodydiode] (S2) at (1,0.5) {};
\draw (1,1) -- (S2.D) (S2.S) -- (1,0);
\node[anchor=west] at ([xshift=0.33cm]S2.center) {$S_{2}$};
% ---- S3 (nmos) ----
\node[nmos,bodydiode] (S3) at (2.35,3.4) {};
\draw (2.35,3.9) -- (S3.D) (S3.S) -- (2.35,2.9);
\node[anchor=west] at ([xshift=0.33cm]S3.center) {$S_{3}$};
% ---- S4 (nmos) ----
\node[nmos,bodydiode] (S4) at (2.35,0.5) {};
\draw (2.35,1) -- (S4.D) (S4.S) -- (2.35,0);
\node[anchor=west] at ([xshift=0.33cm]S4.center) {$S_{4}$};
% ---- Cr (C) ----
\draw (2.6,2.5) to[C, l^=$C_r$] (3.35,2.5);
% ---- Lr (L) ----
\draw (3.35,2.5) to[L, l^=$L_r$] (4.5,2.5);
% ---- Lm (L) ----
\draw (4.5,2.5) to[L, l_=$L_m$] (4.5,1.2);
% ---- T (transformer) ----
\draw (5.15,2.5) to[L] (5.15,1.2);
\draw (5.85,1.2) to[L] (5.85,2.5);
\draw (5.4643,1.2) -- (5.4643,2.5) (5.5357,1.2) -- (5.5357,2.5);
\fill (5.0286,2.0833) circle (1.2pt);
\fill (5.9714,2.0833) circle (1.2pt);
\node at (5.5,2.6917) {$1:n$};
% ---- D1 (D) ----
\draw (6.55,2.5) to[D, l_=$D_1$] (6.55,3.9);
% ---- D2 (D) ----
\draw (6.55,1.5) to[D, l_=$D_2$] (6.55,2.5);
% ---- D3 (D) ----
\draw (7.6,1.2) to[D, l_=$D_3$] (7.6,2.3);
% ---- D4 (D) ----
\draw (7.6,0) to[D, l_=$D_4$] (7.6,1.2);
% ---- Co (C) ----
\draw (8.65,3.9) to[C, l^=$C_o$] (8.65,0);
% ---- RL (R) ----
\draw (9.7,3.9) to[R, l^=$R_L$] (9.7,0);
% ---- annotations ----
\node at (10.65,2.37) {$+$};
\node at (10.65,1.95) {$v_o$};
\node at (10.65,1.53) {$-$};
% ---- crossings without connection ----
\draw (2.2643,2.5) arc[start angle=180,end angle=0,radius=0.12cm];
\draw (6.4643,1.2) arc[start angle=180,end angle=0,radius=0.12cm];
% ---- junction dots ----
\foreach \p in {(1,0),(1,2.5),(1,3.9),(2.35,1.2),(4.5,1.2),(4.5,2.5),(6.55,2.5),(7.6,0),(7.6,1.2),(7.6,3.9),(8.65,0),(8.65,3.9)} \fill \p circle (1.3pt);
\end{circuitikz}
\end{document}
